\documentclass[handout]{beamer}

%%%%%%%%%%%%%%%%%%  PACKAGES %%%%%%%%%%%%%%%%%%
\usepackage{etex}
\usepackage{enumerate}
\usepackage{graphicx}
\usepackage{amsmath}
\usepackage{amssymb}
\usepackage{verbatim}
\usepackage{ifpdf}
\usepackage{bm}
\usepackage{ctable}
\usepackage{epstopdf}
\usepackage{graphicx}
\usepackage{graphics}
\usepackage[english]{babel}
\usepackage{bbm}
%\usetheme{Copenhagen}
%\usetheme{Singapore}
%\usetheme{Malmoe}
\usetheme{CambridgeUS}
%\usetheme{Amsterdam}
\usefonttheme{professionalfonts} % using non standard fonts for beamer
\usecolortheme{beaver}
%\usefonttheme{serif} % default family is serif
\usepackage{bbding}
\usepackage{color}
\usepackage{soul}
\usepackage{isaacshortcuts}
\usepackage{tikz}
\usetikzlibrary{decorations.pathreplacing}
\usetikzlibrary{trees}
%\DeclareMathSizes{5}{5}{5}{5}
%%%%%%%%%%%%%%%%%%  SHORTCUTS STUFF %%%%%%%%%%%%%%%%%%

\def\bi{\begin{itemize}}
\def\ei{\end{itemize}}
\def\ben{\begin{enumerate}}
\def\een{\end{enumerate}}
\def\i{\item [$\color{black}\bullet$]}
\def\bee{\begin{equation*}}
\def\eee{\end{equation*}}
\def\ba{\begin{eqnarray}}
\def\ea{\end{eqnarray}}
\def\baa{\begin{eqnarray*}}
\def\eaa{\end{eqnarray*}}
\def\bg{\beta}
\def\ag{\alpha}
\def\og{\omega}
\def\veg{\varepsilon}
\def\lf{\left}
\def\rg{\right}
\def\mc{\mathcal}
\def\td{\tilde}
\def\v{\vspace{0.2cm}}
\def\vv{\vspace{0.4cm}}




%%%%%%%%%%%%%%%%%%  SLIDES FORMATING %%%%%%%%%%%%%%%%%%
\setbeamertemplate{navigation symbols}{}
\setbeamertemplate{frametitle}[default][center]
\setbeamertemplate{footline}{%
\leavevmode%
\hbox{

%\begin{beamercolorbox}[wd=.25\paperwidth,ht=2.5ex,dp=1ex,center]{date in head/foot}%
%\usebeamerfont{date in head/foot}\insertshortdate
%\end{beamercolorbox}%
\begin{beamercolorbox}[wd=1\paperwidth,ht=2.5ex,dp=1ex,center]{title in head/foot}%
\usebeamerfont{title in head/foot}{F303 Intermediate Investments -  Isaac Hacamo}
\end{beamercolorbox}%
\begin{beamercolorbox}[wd=.25\paperwidth,ht=2.5ex,dp=1ex,center]{pages in head/foot}%
\usebeamerfont{pages in head/foot} \insertframenumber{} / \inserttotalframenumber\hspace*{2ex}
\end{beamercolorbox}

}
\vskip0pt%
}


\begin{document}

\title{\textbf{F303 Intermediate Investments}}
\author{Lecture 14: Interest Rate Risk and Duration \vspace{-1cm}}
\date{}
\frame[plain]{\titlepage}
\setcounter{framenumber}{0}


%%%%%%%%%%%%%%%%%%  SLIDE  %%%%%%%%%%%%%%%%%%
\frame{\frametitle{\large{\textbf{Goals for Today}}}
\bi
\i Interest Rate Risk
\vv
\i Duration and Immunization
\vv
\i Problem Solving
\ei
}



\frame{\frametitle{\large{\textbf{Coupon Rate and Interest Rate Sensitivity}}}
\begin{center}
Example 1: Prices of 8\% Coupon Bond (semi-annual payments)
\bt{|c|c|c|c|}
\hline
Yield to Maturity & T = 1 Year & T = 10 Years & T = 20 Years\\
\hline
8\% & 1,000 & 1,000 & 1,000 \\
\hline
9\% & 990.64 & 934.96 & 907.99 \\
\hline
Change in Price & \color{red} 0.94\% & \color{red}6.50\% & \color{red}9.20\%\\
\hline
\et
\end{center}

\begin{center}
Example 2: Prices of Zero-Coupon Bond (semi-annual compounding)
\bt{|c|c|c|c|}
\hline
Yield to Maturity & T = 1 Year & T = 10 Years & T = 20 Years \\
\hline
8\% & 924.56 & 456.39 & 208.29\\
\hline
9\% & 915.73 & 414.64 & 171.93\\
\hline
Change in Price & \color{red}  0.96\% & \color{red} 9.15\% & \color{red} 17.46\%\\
\hline
\et
\end{center}
}


\frame{\frametitle{\large{\textbf{Review from Last Lecture: Interest Rate Sensitivity}}}
\bi
\i Prices of \textbf{long-term bonds} tend to be more sensitive to interest rate changes than prices of short-term bonds; 
\vv
\i Interest rate risk is inversely related to the \textbf{bond's coupon rate} (low coupon bonds are more sensitive)
\vv
\i As \textbf{maturity increases}, price sensitivity to yield changes increases at a decreasing rate;
\vv
 \i  Both maturity and the level of the coupons matter for interest rate sensitivity. We need to put them together. We need a better measure: \textbf{a measure of effective maturity!}
\ei
}



\frame{\frametitle{\large{\textbf{Duration}}}
\bi
\i To understand the interest rate sensitivity of a bond we need to know \ul{when we receive the cashflows}, and \ul{what it the size of the cash flows.} In other words, the timing and the amount of cash flows paid by the bond is important for interest rate sensitivity. 
\vv
\i \textbf{Concept of duration}: The weighted average of the times at each coupon or principal payment is made by the bond. It provides an 'weighted average' timing when you will receive your cash flows. Even though you will receive your bond payments spread over time, duration of a bond is providing \textbf{one hypothetical point in time} when you will receive your cash-flows.
\ei
}
\frame{\frametitle{\large{\textbf{Duration: Zero-coupon bonds}}}
\bi
\i According to the definition from the previous page, what is the duration of a:
\vv
\bi
\i one year zero-coupon bond? 
\i two-year zero coupon bond? 
\i three-year zero coupon bond? 
\ei
\vv
For a zero-coupon does the duration coincide with the actual payment timing?   
\ei
}
%
%\frame{\frametitle{\large{\textbf{Duration: Zero-coupon bonds}}}
%\bi
%\i According to the definition from the previous page, what is the duration of a:
%\vv
%\bi
%\i one year zero-coupon bond? \color{red} D=1 \color{black}
%\i two-year zero coupon bond? \color{red} D=2 \color{black}
%\i three-year zero coupon bond? \color{red} D=3 \color{black}
%\ei
%\vv
%For a zero-coupon does the duration coincide with the actual payment timing?   \color{red} Yes, it does.\color{black}
%\ei
%}

\frame{\frametitle{\large{\textbf{Duration: Coupon bonds}}}
\bi
\i Now that we know the duration of a zero coupon bond, let's try to figure the \textbf{duration of a three coupon bond}. The three year coupon bond pays three coupons and the maturity value. \ul{Let's think of each payment as if it is a zero coupon bond.} What is the duration of each payment in this context? \\
\begin{center}
\begin{tabular}{ccccl}
t= 0 & t=1 & t=2 & t=3 & \\
\hline
\\
 & $\frac{C}{(1+y_1)}$ &   &   & $D_1$= \color{red} 1\color{black}\\

 \hline
 \\
  &  & $\frac{C}{(1+y_2)^2}$    &   & $D_2$= \color{red} 2\color{black}\\ \\
  \hline
 \\
& &   & $\frac{C+F}{(1+y_3)^3}$ &    $D_3$= \color{red} 3\color{black}\\ \\
\end{tabular}
\end{center}
\ei
}

\frame{\frametitle{\large{\textbf{Duration: Coupon bonds}}}
\bi
\i The duration of the coupon bond ($D_{c}$) should be some sort of weighted average of the durations of each payment. Using the values from the previous table, we could take the \ul{weighted average of each duration using as weight the relative size of the payments in each period}, that is:
\baa
D_{c}=\underbrace{D_1}_{=1} \frac{C/(1+y_1)}{Price}+\underbrace{D_2}_{=2}\frac{C/(1+y_2)^2}{Price}+\underbrace{D_3}_{=3} \frac{(F+C)/(1+y_3)^3}{Price}
\eaa
Not that the sum of all weights is equal to the price of the bond, that is:
\baa
Price=C/(1+y_1)+C/(1+y_2)^2+(F+C)/(1+y_3)^3
\eaa
\ei
}


\frame{\frametitle{\large{\textbf{Duration: General Formula}}}
The weight, $w$, associated with the cash flow ($CF$) made at time $t$ is given by:
\baa
w_t=\frac{CF_t/(1+y_t)^t}{Price}
\eaa
Then the Macaulay's duration is defined as:
\baa
D=\sum_{t=1}^Tt\times w_t
\eaa
\baa
D=1\times \frac{C/(1+y_1)}{Price}+2\times \frac{C/(1+y_2)^2}{Price}+\dots +T\times \frac{(F+C)/(1+y_T)^T}{Price}
\eaa
}


\frame{\frametitle{\large{\textbf{Fulcrum Point and Duration}}}

\begin{center}
\ul{You can think of duration as a fulcrum point.}

\vspace{0.5cm}

\includegraphics[scale=0.5]{./duration2}
\end{center}
}

\frame{\frametitle{\large{\textbf{Duration}}}
\bi
\i The cash flow weights are \textbf{related to the importance of every single payment} to the value of the bond. 
\vv
\i The weight applied to each payment time is the proportion of the total value of the bond accounted by that payment. 
\vv
\i The proportion is the present value of the payment divided by the bond price.
\ei
}


\frame[t]{\frametitle{\large{\textbf{Practice Problem: Calculating Duration}}}
Compute the duration of a 2-year 8\% coupon bond with semi-annual payments. Assume YTM=5\% and the face value is \$100. 
}

\frame[t]{\frametitle{\large{\textbf{Practice Problem: Calculating Duration}}}
Compute the duration of a 2-year 8\% coupon bond with semi-annual payments. Assume YTM=5\% and the face value is \$100. 
\color{red} The price of this bond is 105.64 - use your calculator to compute this value, remember that this is a semi-annual bond, as a result, you should divide both the YTM and coupon by two. The YTM and coupon rate are both APRs, not EARs!
\baa
D&=&0.5\times\frac{4/105.64}{1+0.025}+1\times\frac{4/105.64}{(1+0.025)^{2}} \\
&+&1.5\times\frac{4/105.64}{(1+0.025)^{3}}+2\times\frac{104/105.64}{(1+0.025)^{4}}\\
&=&1.89
\eaa
The duration is close to the maturity value because the maturity value is very large relative to the coupon paymenta, thus, the weight on the last payment period $(t=2)$ is very large.
}



\frame{\frametitle{\large{\textbf{Duration - Its Importance}}}
\bi
\i Duration is a key concept in fixed income portfolio management. It is a simple summary statistic of the effective average maturity of the portfolio;
\vv
\i  It is a measure of the interest rate sensitivity of the portfolio, and it is an essential tool in immunizing portfolio from interest rate risk;
\vv
\i  Long-term bonds are more sensitive to interest rate movements than are short-term bonds. The duration measure enables us to quantify this relationship.
\vv
\i When interest rates change, the proportional change in a bond's price can be related to the change in its yield to maturity.
\ei
}

\frame{\frametitle{\large{\textbf{Bond Price Changes}}}
Bond price volatility is proportional to the bond's duration; therefore duration becomes a natural measure of interest rate exposure.
\baa
\frac{\Delta P}{P}=-D\times \frac{\Delta y}{1+y}
\eaa
where $y$= annual YTM and $D$= duration expressed in years.
}


\frame{\frametitle{\large{\textbf{Duration of a Portfolio}}}
\bi
\i The duration of a portfolio is the portfolio of the duration, meaning:
\baa
D_{P}=\frac{V_{1}}{V}D_{1}+\frac{V_{2}}{V}D_{2}+\dots +\frac{V_{M}}{V}D_{M}
\eaa
where $V_{i}$ is the value of the security $i$, $V$ is the total value of the portfolio, and $D_i$ is the duration of security $i$.
\vv
\i \textbf{Important:} In the above formula, $V_i$ is positive for an asset, and $V_i$ is negative for a liability. Also remember that 
\baa
V=V_1+V_2+\dots + V_M
\eaa
\ei
}


\frame{\frametitle{\large{\textbf{Immunization}}}
\bi
\i Immunization techniques refer to strategies  used by investors to shield their overall financial position from exposure to interest rate fluctuations. 
\vv
\i If this is not taken into consideration, a change  in interest rates could run havoc with the institution's balance sheet, threatening bankruptcy.
\ei

}

\frame{\frametitle{\large{\textbf{Offsetting Risks}}}
\bi
\i Fixed income investors face two offsetting types of interest rate risks:
\vv
\bi
\i \textbf{Price risk}; and
\i \textbf{Reinvestment rate risk}.
\ei
\vv
\i Increases in interest rates cause capital losses but at the same time increase the rate at which reinvested income will grow.
\ei
}

\frame{\frametitle{\large{\textbf{So How Would Duration Help?}}}

\bi
\i If the portfolio \textbf{duration is chosen appropriately}, these \textbf{two effects will cancel out} exactly. 
\vv
\i When the portfolio's duration is set equal to the  investor's horizon date, the accumulated value  of the investment fund at the horizon date will be  unaffected by interest rate fluctuations.
\ei
}



\frame[t]{\frametitle{\large{\textbf{Problem 1}}}
An investor has \$100,000 to invest. She borrows another \$80,000 by issuing 10-year zero coupon bonds worth \$80,000 today. She invests \$70,000 in 3-year zero coupon bonds and \$110,000 in a coupon bond with duration of 4 years. \\
\vv
a) What is the duration of her portfolio?\\

}

\frame[t]{\frametitle{\large{\textbf{Problem 1}}}
An investor has \$100,000 to invest. She borrows another \$80,000 by issuing 10-year zero coupon bonds worth \$80,000 today. She invests \$70,000 in 3-year zero coupon bonds and \$110,000 in a coupon bond with duration of 4 years. \\
\vv
a) What is the duration of her portfolio?\\

\color{red}
The total value of her portfolio is \$100,000 (=110k+70k-80k).
\baa
D_p&=&\frac{70}{100}\times 3 +\frac{110}{100}\times 4 -\frac{80}{100}\times 10=-1.5
\eaa
}

\frame[t]{\frametitle{\large{\textbf{Problem 1 (cont)}}}
b) The yield curve is currently flat at 4\%. If the yield curve shifts up by 0.5\% in a few days, will she have a gain or a loss? Calculate it in percent.
}

\frame[t]{\frametitle{\large{\textbf{Problem 1 (cont)}}}
b) The yield curve is currently flat at 4\%. If the yield curve shifts up by 0.5\% in a few days, will she have a gain or a loss? Calculate it in percent.
\color{red}
\baa
\frac{\Delta p}{p}=-\frac{(-1.5)}{1+0.04}\times 0.005=0.00721=0.721\%
\eaa

}

\frame[t]{\frametitle{\large{\textbf{Problem 1 (cont)}}}
c) Assume that she can choose the duration of the coupon bond. What bond duration should she pick to have a portfolio with no interest rate risk?\\
}

\frame[t]{\frametitle{\large{\textbf{Problem 1 (cont)}}}
c) Assume that she can choose the duration of the coupon bond. What bond duration should she pick to have a portfolio with no interest rate risk?\\
\color{red}
A bond has no interest rate risk if duration is equal to zero. 
\baa
D_p&=&\frac{70}{100}\times 3 +\frac{110}{100}\times D_c -\frac{80}{100}\times 10=0
\\
\rightarrow D_c&=&5.36\ \text{years}
\eaa

}

\frame[t]{\frametitle{\large{\textbf{Problem 1 (cont)}}}	
d) Assume that she can borrow any amount, and that she wants to always invest the same amount in the 3-year zero coupon bond and in the coupon bond with duration of 4 years. How much should she borrow to have a portfolio with no interest rate risk?\\
}

\frame[t]{\frametitle{\large{\textbf{Problem 1 (cont)}}}
d) Assume that she can borrow any amount, and that she wants to always invest the same amount in the 3-year zero coupon bond and in the coupon bond with duration of 4 years. How much should she borrow to have a portfolio with no interest rate risk?\\
\color{red}

Here we also want to impose that the duration of the portfolio is  zero, but in this case we want to choose the amount to borrow, while imposing that the amount invested in the other two assets is always equal.
\baa
D=0&=&\frac{50+x/2}{100}\times 3 +\frac{50+x/2}{100}\times 4 -\frac{x}{100}\times 10\\
0&=&(50+x/2)\times 3 +(50+x/2)\times 4 -x\times 10\\
0&=&   350 -x\times 13/2\\
\rightarrow x&=&53,846.15
\eaa
}



%%%%%%%%%%%%%%%%%%  SLIDE  %%%%%%%%%%%%%%%%%%

\newcounter{finalframe}
\setcounter{finalframe}{\value{framenumber}}

\setcounter{framenumber}{\value{finalframe}}
\end{document}


%%%%%%%%%%%%%%%%%  SLIDE  %%%%%%%%%%%%%%%%%%






